\documentclass[10pt]{article}
\usepackage[margin=20mm]{geometry}
\usepackage{amsmath,amssymb,mathtools,fontspec,unicode-math,xurl,hyperref}
\IfFontExistsTF{DejaVu Serif}{\setmainfont{DejaVu Serif}}{\setmainfont{Latin Modern Roman}}
\setmathfont{latinmodern-math.otf}
\setlength{\parindent}{0pt}\setlength{\parskip}{6pt}\sloppy
\title{Technical Verification of
Exact Enumeration of Fullerenes}
\author{}
\date{26 October 2026}
\begin{document}
\maketitle
Abstract: The explicit formula in Theorem 3.26 does not have the nonnegative integral coefficients required of the stated counting function. This remains true after restoring an evidently malformed exponent. Two further coefficient changes remove the obstruction and agree with every entry of the supplied finite enumeration table. An independent construction verifies the claimed dimension and finite-coefficient injectivity of the ambient space. Five local corrections concern the general Weil-level statement, stratum and marking distinctions, and normalization or arithmetic displays. The noncompact weight-two theta integral requires a specified regularization. These findings require correction of the published formula, but do not establish failure of the underlying enumeration strategy. This is an initial technical review; scientific adjudication remains pending.\par
\section*{Paper identity, conventions, and review boundary}
This report concerns only M25-0150: Philip Engel and Peter Smillie, Exact enumeration of fullerenes, with an appendix by Jan Goedgebeur, Duke Mathematical Journal 174(3) (2025), 575–613, DOI 10.1215/00127094-2024-0037 (https://doi.org/10.1215/00127094-2024-0037). The controlling input is the 39-page journal-typeset PDF deposited at NSF PAR, record 10642973. Journal page \(p\) is PDF page \(p-574\). The separately acquired arXiv:2304.01655v3 PDF and its TeX source were used for version comparison. The findings below are confined to these hash-pinned inputs, not to an unexamined later correction.\par
All 39 journal pages were read in extracted text. Relevant formula, proof, table, and dependency pages were also inspected as rendered images. The complete 14-file input manifest was parsed and its byte lengths and SHA256 hashes checked. Native session auxiliaries performed bounded geometry, enumeration, and modular-form checks; their actual roles and separate records are preserved. They were not independent human referees or independently verified model instances.\par
The paper counts oriented fullerenes with \(2n\) vertices: mirror images are distinct unless related by an orientation-preserving equivalence. Euler’s formula gives twelve pentagons and \(n-10\) hexagons, hence \(\mathrm{BB}_n=0\) for \(n<10\). In the dual triangulation \(q^n\) records \(2n\) triangles. We use the paper’s normalization \[E_2(q)=-\frac1{24}+\sum_{n\ge1}\sigma_1(n)q^n,\qquad
 E(q)=\frac16+\sum_{n\ge1}\left(\sum_{d\mid n}\chi_3(d)\right)q^n,\] where \(\chi_3(1)=1,\ \chi_3(2)=-1,\ \chi_3(0)=0\). For higher Eisenstein series, the divisor sums and generalized Bernoulli constants were checked against Definition 3.25 and the primary normalization reference [Retained latex: \textbackslash{}cite[Section 5]\{Stein\}]. Positive-degree counterexamples below do not depend on a choice of additive constant.\par
\section*{Problems with decisive evidence}
Problem 1 (Major: the main explicit generating series).  Location. Theorem 3.26, journal pp. 605–606 (PDF pp. 31–32); arXiv v3 PDF p. 25; pinned TeX lines 1672, 1689–1690.\par
Evidence. The denominator of the \(E_4(q^3)\) term is typeset with only the first digit of \(11\) in the exponent of \(2\). The source is literally `2\textasciicircum{}11`, rather than `2\textasciicircum{}\{11\}`. Read literally, its denominator is \(2\cdot1\cdot3^9\cdot5\), and the displayed series has \[[q^3]\mathrm{BB}_{\rm literal}(q)=-\frac{911198717}{67184640}.\] The source’s commented coefficient vector instead gives denominator \(201553920=2^{11}3^9 5\), supporting restoration of the missing braces.\par
Equation reference: eq:bad. Even after making this favorable restoration, the formula still fails. Let \(F\) denote that exponent-restored series, with all other printed coefficients unchanged, and put \[R=F-\frac1{64}E_2(q^9)-\frac2{27}E_2(q^{12}).\] Exact rational arithmetic applied to the other 43 terms gives \[[q^9]R=-\frac1{18},\qquad [q^{12}]R=\frac{47}{48}.\] Consequently, \[
 [q^9]F=-\frac1{18}+\frac1{64}=-\frac{23}{576},
 \qquad
 [q^{12}]F=\frac{47}{48}+\frac2{27}=\frac{455}{432}.\] The first violates the elementary vanishing condition \(n<10\); the second is not an integer. These are decisive failures of the stated theorem, irrespective of the correctness of the appendix’s numerical counts. Two separately written reviewer calculations, based on independent transcriptions, agree. The full per-term sums are retained.\par
Repair and its boundary. In addition to restoring \(2^{11}\), replace \[\frac1{64}E_2(q^9)+\frac2{27}E_2(q^{12})
 \quad\text{by}\quad
 \frac1{18}E_2(q^9)+\frac1{48}E_2(q^{12}).\] The resulting series \(F^\ast\) agrees exactly with all 201 supplied table entries \(0\le n\le200\), including its zero constant term. Before these two changes, the exponent-restored formula differs at 34 of these indices. This comparison uses the published table; it is not a fresh generation of all fullerene graphs.\par
The difference is \[F^\ast-F=\frac{23}{1728}\bigl(3E_2(q^9)-4E_2(q^{12})\bigr).\] The parenthesis equals one third of \[\bigl(E_2(q)-12E_2(q^{12})\bigr)
 -\bigl(E_2(q)-9E_2(q^9)\bigr),\] so it is a weight-two modular form on \(\Gamma_0(36)\), contained in the claimed ambient space \(\Omega\). Section 4 (\#sec:positive) independently verifies coefficient injectivity on \(\Omega\). Thus, conditional on the inclusion \(\mathrm{BB}\in\Omega\) and the accuracy of the supplied initial counts, the repaired candidate is the unique generating series in that space with those counts. No author adoption of these corrections is claimed.\par
Counterevidence limiting the impact. The embedded commented Sage appendix contains coefficient vectors compatible with the proposed changes; its original Sage basis ordering was not regenerated. The correction is supported only at indices divisible by \(9\) or \(12\), apart from the constant. Example 3.27’s restricted polynomial and Example 3.28’s count \[\mathrm{BB}_{5000}=1203397779055806181762759\] survive the two weight-two changes. Their survival does not repair the formula for general \(n\). The severity is Major because an explicit principal theorem gives impossible counts, even though a short correction appears sufficient.\par
Problem 2 (Minor: exponent of the discriminant group is not the Weil level).  Location. Proposition 3.9 and its explanation, journal p. 598 (PDF p. 24).\par
The general assertion for an even lattice of even rank substitutes \(N\)-torsion of its discriminant group for level of its quadratic discriminant form. Take \(L=A_1\oplus A_1\), with Gram matrix \(2I_2\). Its discriminant group has exponent two, but for \(\gamma=(1/2,0)\), \[\theta_\gamma(\tau)=\sum_{m,n\in\mathbb Z}q^{(m+1/2)^2+n^2}
 =2q^{1/4}+\cdots,\qquad
 \theta_\gamma(\tau+2)=-\theta_\gamma(\tau).\] Thus \(T^2\in\Gamma(2)\) acts nontrivially. The correct level is the least \(N\) such that \(N(\gamma,\gamma)/2\in\mathbb Z\) for every discriminant class; here it is four. This is also the condition in [Borcherds, p. 13] and [Retained latex: \textbackslash{}cite[p.\textasciitilde{}5, Propositions 4.3--4.5]\{Scheithauer\}].\par
The counterexample invalidates the unrestricted formulation, not the needed levels for the actual Eisenstein lattices. Independent Gram-matrix calculations give quadratic levels \(3,2,3,1\) for \(A_2,D_4,E_6,E_8\). Taking negatives and isotropic subquotients does not enlarge their least-common-multiple level. Keeping the quadratic forms in the complement and saturation argument therefore supports the levels \(1,2,3,6\) used in Theorems 3.13 and 3.16. Replace the general hypothesis and make this application explicit; enlarging every asserted level merely because the discriminant group was misdescribed is not warranted.\par
There is a related, readily repairable hypothesis mismatch in Proposition 3.18 (p. 602) and its use in Theorem 3.20 (p. 603): the former assumes \(d\mid N\), although some applications take \(d=|G|\) without that divisibility. The same argument works for every positive integer \(d\). If \(h\) is modular for \(\Gamma(N)\) and has a \(q^{1/d}\)-expansion, then it is invariant under \(T^d\). For \[M=\begin{pmatrix}1+aN&bd\\cN&1+eN\end{pmatrix},
 \qquad \det M=1,\] one has \(M=T^{bd}K\) with \(K\in\Gamma(N)\). Conjugation by \(\operatorname{diag}(d,1)\) yields \(\left(\begin{smallmatrix}1+aN&b\\cdN&1+eN\end{smallmatrix}\right)\in\Gamma(N,d)\). Removing the unnecessary divisibility assumption repairs this invocation; it gives no counterexample to the final level bound.\par
Problem 3 (Minor: the stratum description needs action and collision data).  Location. Proposition 2.14 and the criterion after Definition 2.15, p. 585; Proposition 2.19, pp. 587–588. Relevant corrective evidence occurs in Table 2, pp. 590–593, and Example 3.30, pp. 607–608.\par
Conjugacy of inertia is finer than the abstract automorphism group together with the upstairs curvature multiset. Consider \(G=C_2\) and upstairs profile \((2,2,1^8)\). Table 2 distinguishes quotient profiles \[(3^e,3^e;2,1,1,1,1),\qquad (4,4;1,1,1,1),\] of dimensions four and three. Generic cone configurations invariant under \(z\mapsto-z\) realize the first with curvature-two points at \(\pm1\), and the second with curvature-two points at \(0,\infty\); the eight remaining curvature-one points form four generic pairs. Both have full automorphism group \(C_2\). In the first case the involution exchanges the two \(C_3\) collision factors, producing nonabelian inertia. In the second it fixes the collision lines and centralizes their order-three reflections; the central kernel has cyclic quotient, hence the inertia is abelian. The two inertia types cannot be conjugate. The action on collision data, not just a jump in group order, matters.\par
The numerical criterion for a nonempty exact bistratum also needs exclusion of forced extra symmetry. For \(G=C_1\) and profile \((4,4,4)\), the three equal cone angles give the doubled equilateral triangle with full \(D_3\) symmetry. The later table explicitly marks the repetition. Finally, Proposition 2.19’s bijection with unmarked completed triangulations is too coarse: Example 3.30 exhibits two collision branches of the same \((5,5,2)\) triangulation, according as the original curvature-two point remains separate or joins a curvature-five cluster. Thurston’s completion retains such partitions [Thurston, p. 528].\par
These distinctions should be included in the statement: use the ramification action and the normalization with marked collision data, then sum inverse stabilizer orders over its finite fibres. The quotient by \(N_H(G)/G\) handles residual automorphisms of an actual \(G\)-quotient. The paper’s detailed table and worked example already retain the needed refinements, and the resulting conversion remains rational and triangular with the appropriate exact-stratum diagonal. No counterexample to that refined inversion has been established. A complete \(96\times96\) transition matrix was not independently computed.\par
Problem 4 (Minor: the exponent unit in a zero-dimensional stratum).  Location. Example 2.22, p. 589 (PDF p. 15), compared with Definition 2.21 and p. 604.\par
The parameter \(r\) in \(E(q^r)\) must be half the minimum number of triangles, since \(q^n\) records \(2n\) triangles. The doubled equilateral triangle has two triangles and gives \(E(q)\), not \(E(q^2)\). Likewise the icosahedral quotient has \(20/60=1/3\) triangle, hence \(r=1/6\); pullback gives \(E(q^{10})\), as required. The table entries and the later statement on p. 604 already use the correct convention. Change the explanatory sentence; no propagated error in the table is shown.\par
Problem 5 (Minor: the sign of the constant coefficient).  Location. Remark 3.14, p. 601 (PDF p. 27).\par
For \(d=\operatorname{len}(\kappa)-3\), the displayed positive volume formula lacks \((-1)^d\). Indeed, for \(y=\operatorname{Im}\tau\) and \(D=(2\pi i)^{-1}\partial_\tau\), \[D^d y^d=\frac{(-1)^d d!}{(4\pi)^d}.\] In the convergent case the zero-vector term therefore gives \[c_0=\frac{(-1)^d d!\mathop{\mathrm{Vol}}(P\Gamma_\kappa\backslash\mathbb{CH}^d)}
               {6(4\pi)^d}.\] A concrete check is \(\kappa=(1^{12})\), \(d=9\). The weighted series is \(C E_{10}\), where \[C=\frac{809}{2^{15}3^{13}25}>0,\qquad
 [q^0]CE_{10}=-\frac{809}{344802322022400}<0.\] This contradicts the printed positive-volume expression. The derivative sign and the corresponding volume formula agree with [Retained latex: \textbackslash{}cite[Corollary 4.14, p.\textasciitilde{}19]\{ES\}]. This is a constant-term normalization correction; no additional failure of positive-degree enumeration follows. A cusp-regularized weight-two constant requires the separate specification discussed below.\par
Problem 6 (Minor: one intermediate stabilizer-mass denominator).  Location. Example 3.30, p. 608 (PDF p. 34), in the sum immediately following the correct individual \(T_2\) weight.\par
The individual contribution is \(1/(2^{11}3^{10}5)\), but the subsequent sum uses \(2^{12}\). With common denominator \(2^{15}3^{10}25\), the correct three numerators are \[1+80+3200=3281=17\cdot193;\] the displayed intermediate denominator instead gives \(1+40+3200=3241\). Thus the final numerator on that page is correct and the middle denominator should be \(2^{11}\). The factor \(5\) is present in the actual image and is not an error.\par
Further counterevidence is that the seven \(q^2\) masses sum to \(697/40310784\), and both initial coefficients agree with the final \(A=1/(2^{16}3^{10}25)\), \(B=1/(2^{16}3^2 25)\) on p. 610. The error is confined to an intermediate display; it does not refute those final weight-nine coefficients.\par
\section*{Analytic evidence still needed}
Problem 7 (Needs Evidence: regularization at a noncompact weight-two cusp).  Location. Definition 3.7 and the sketch of Theorem 3.8, p. 597 (PDF p. 23). The definition gives an ordinary theta integral also for signature \((1,1)\); the sketch calls its noncompact case conditionally convergent.\par
There is a direct obstruction to that literal convergence assertion. Put \(\mathcal O=\mathbb Z[\zeta_6]\), \(\alpha=1+\zeta_6\), and take the Hermitian lattice \[\Lambda=\mathcal Oe+\mathcal Of,\qquad
 \star=\begin{pmatrix}0&\alpha\\\overline\alpha&0\end{pmatrix}.\] Its underlying form \((2/3)\operatorname{Re}\star\) is even, unimodular, and of signature \((2,2)\): in the indicated Eisenstein basis the off-diagonal integral block is \(\left(\begin{smallmatrix}1&1\\0&1\end{smallmatrix}\right)\). It has an isotropic sublattice \(\mathcal Oe\). At \(\tau=iy\), all terms of the theta kernel are positive. At cusp height \(H\), the terms indexed by \(\beta e\) have the form \[\exp(-Cy|\beta|^2/H),\qquad C>0.\] There are at least \(cH\) Eisenstein integers with \(|\beta|\le\sqrt H\), so the kernel is bounded below by \(c_yH\). Integration over a cusp strip with invariant measure proportional to \(dx\,dH/H^2\) therefore dominates \[c'_y\int_{H_0}^{\infty}\frac{dH}{H}=\infty.\] A positive divergent integral is not conditionally convergent. This boundary occurs in the paper’s \((5,5,1,1)\) profile: its \(E_8\oplus E_8\) root complement in \(\mathrm{II}_{2,18}\) leaves \(\mathrm{II}_{2,2}\).\par
The cited earlier construction explicitly assumes \(s>1\) in [Retained latex: \textbackslash{}cite[Definition 4.4, p.\textasciitilde{}13]\{ES\}]. Its isotropic sum in Lemma 4.10 is proportional to \(\sum_{\beta\ne0}|\beta|^{-2s}\), which diverges at \(s=1\). Thus that reference does not directly supply the missing ordinary integral.\par
Required evidence and limit. Specify a truncation or regularized lift, its subtraction convention, and the derivative or holomorphic-part operation; justify both the transformed weight-two completion and the stated positive Fourier coefficients. The leading divergent term after the \(y\)-prefactor can be independent of \(\tau\), so differentiation may remove it in a correctly defined construction. The paper itself recognizes the exceptional case and a nonholomorphic correction. No counterexample to the resulting quasimodularity is asserted here. The demonstrated defect is in the unregularized route; the status of the required regularized argument is Needs Evidence.\par
\section*{Affirmative checks and core dependencies}
\subsection*{Independent dimension and coefficient injectivity}
Propositions 3.22–3.23 were checked beyond repeating their software claim. Let \(\chi_3\) be as above. The nine modular summands of \(\Omega\), after identifying the \(-I\) cosets of the paper’s groups, have the following descriptions: \[\begin{array}{c|ccccccccc}
\text{weight}&10&9&8&7&6&5&4&3&2\\ \hline
\text{level}&1&3&6&6&12&12&72&72&72\\
\text{character}&1&\chi_3&1&\chi_3&1&\chi_3&1&\chi_3&1\\
\text{dimension}&1&4&9&8&13&11&44&32&20
\end{array}\] The level row denotes \(\Gamma_0(N)\) with the displayed character. The dimensions sum to 142. The two additional \(E_2\) series, seven dilates \(E(q^t)\), \(t\in\{1,2,3,4,5,6,10\}\), and a constant give an upper bound of 152.\par
For \(N=6,12,72\), the torsion-free modular curves have respectively \[(\mu,c,g)=(12,4,0),\ (24,6,0),\ (144,16,5).\] The character is trivial on the primitive cusp parabolics. The weight-one character line bundle \(L\) has \(L^2=K(D_{\rm cusp})\), hence degrees \(1,2,12\). In all weights used here, Riemann–Roch gives \(\dim H^0(L^k)=k\deg L+1-g\). At level three, the weight-three character bundle has trivial elliptic stabilizer and degree one on a genus-zero curve, so its cube has four sections. The weight-ten level-one dimension is one by the valence bound. These are explicit upper bounds independent of a software dimension command.\par
For lower bounds and injectivity, reviewer code constructed integral Fourier series \[\Theta_t=1+6\sum_{n\ge1}\sum_{a\mid n}\chi_3(a)q^{tn},
 \qquad
 B=\sum_{n\ge1}\sum_{a\mid n}\chi_3(n/a)a^2q^n.\] Here \(\Theta_t\) is a dilated \(A_2\) theta series of weight one and \(B\) a weight-three level-three Eisenstein series. Homogeneous products with \(3t\mid N\), together with \(\Theta_1^{9-3j}B^j\) for \(0\le j\le3\), reach the nine upper bounds. The correct weight, character and level of every candidate were checked. Adjoining the ten extra series gives a \(200\times152\) coefficient matrix for \(q^0,\ldots,q^{199}\).\par
Over the verified prime \(p=1000003\), a selected \(152\times152\) minor has determinant \[\det A_{\rm minor}\equiv859210\pmod{1000003}.\] All columns lift to integral series, so the determinant is also nonzero over \(\mathbb Q\) and \(\mathbb C\). A separately written root-executor row elimination checked this saved minor directly. The largest selected coefficient index is 168. This establishes both dimension 152 and injectivity at the stated precision without using the printed answer or fullerene counts. The appendix’s \(q^{200}\) value is an additional check, not an off-by-one failure.\par
\subsection*{Geometry, normalization, and finite computations}
The bounded quotient-data enumeration produces 116 covering-group labels; removing the 20 repetitions explicitly identified in the table leaves 96. Every label passes the upstairs total-curvature-twelve condition. This checks the finite catalogue arithmetic under its stated classification; it does not construct every stratum or a complete transition matrix.\par
The period-integrality argument, the face and edge period superlattices, the cyclic-to-dihedral enlargement for orders four through six, and the reciprocal-stabilizer interpretation were examined. The positive-vector orbital integral in [Retained latex: \textbackslash{}cite[Lemma 4.8, pp.\textasciitilde{}14--15]\{ES\}] gives the stated inverse-stabilizer coefficient after differentiation. The finite-index and finite-center normalizations were checked. The full vector-valued extension for every discriminant coset was not reproved.\par
Independent rational arithmetic verified the formula witnesses, constant cancellation after repair, the actual root-lattice levels, and the worked mass example. A separate enumeration auxiliary compared all 201 table entries. Explicit 20- and 24-vertex examples and their preserving/reversing automorphism counts were checked by reviewer code; these establish examples and the orientation convention, not classification of all graphs. The restricted polynomial of Example 3.27 agrees with the repaired formula at every admissible index through 200, including the empty \(n=1\) case. The 10,000-vertex value and its stated approximate asymptotic discrepancy also check.\par
The logical boundary is therefore precise: finite-coefficient injectivity is independently established, while the identification of the true counting series with an element of \(\Omega\) still uses the refined geometric counting argument and the analytic input identified in Problem Problem 7 (\#prob:cusp). Agreement with a supplied table alone is not promoted to an unconditional all-\(n\) theorem.\par
\section*{Assessment}
Item | Status | Established impact\par
Problem 1 | Major | Main explicit series gives impossible counts\par
Problems 2–6 | Minor | Local hypotheses, distinctions, or normalizations\par
Problem 7 | Needs Evidence | Regularized weight-two analytic construction\par
Dimension and truncation | Verified in scope | Rank 152 with a valid modular basis\par
All supplied initial counts | Compared | Repaired formula matches \(0\le n\le200\)\par
The explicit principal formula needs correction. Its identified failure is not evidence that every theorem is false or that the method is beyond repair. The table comparison, modular correction, and independent injectivity certificate provide a concrete route to a corrected formula, conditional on the stated counting and analytic dependencies. A corrected exposition should also retain the marked boundary data and specify the weight-two regularization. The review records `scientific\_decision=pending` and `human\_reviews=[]`; neither acceptance nor independent adjudication has been performed.\par
\section*{Supporting checks}
No other paper or controller state was modified by this review. The batch token is \[\texttt{b72d1a0d-a138-4ff7-9d37-932d176027ca}.\] The full scientific prompt and frozen contract 0.7.1 govern the review. The UI selection Personal account / Chat / GPT-6 / Pro was observed by the external controller; this executor did not independently verify the backend.\par
The principal PDF is 643051 bytes, SHA256\par
> `408ef7c8c55b4a74f99d93c7df6b30c0c8903fa6cc601a0f3ac035e5dc2eb850`.\par
The supplied arXiv v3 PDF (772905 bytes) and TeX file (131064 bytes) are separately hash-bound in the complete input manifest. The coefficient-matrix CSV has SHA256\par
> `1b744527024355ee85f5294051369d71051af0fbdde5ea142a82fc496c4e50e3`.\par
The source TeX contains a commented Sage Appendix B, lines 2203–2347, which was read statically. It is therefore incorrect to report the computation’s source as wholly unavailable. No separate version-pinned Sage environment or matching execution log was established. The supplied buckygen archive was identified, but its exact revision used for the appendix’s enumeration was not verified. Neither author program was executed; missing execution provenance is a reproducibility limitation, not a mathematical counterexample.\par
The evidence bundle retains reviewer code and calculation attempts, per-term sums, all table comparisons, dependency sources, actual image views, and individual hashes. The literal and exponent-restored formula readings are distinguished in `exponent\_interpretation.json`. File-inventory parsing is not counted as scientific reading.\par
9 R. E. Borcherds, Automorphic forms with singularities on Grassmannians, author manuscript dated 29 September 1996, corrected 29 May 1997, Section 4, pp. 12–14. https://math.berkeley.edu/\textasciitilde{}reb/papers/aut/aut.pdf.\par
P. Engel and P. Smillie, The number of convex tilings of the sphere by triangles, squares, or hexagons, arXiv:1702.02614v2, 16 January 2018, pp. 12–19. https://arxiv.org/pdf/1702.02614v2.\par
N. R. Scheithauer, The Weil representation of \(\mathrm{SL}_2(\mathbb Z)\) and some applications, hash-pinned author-hosted PDF, pp. 5, 14–15. https://www.mathematik.tu-darmstadt.de/media/algebra/homepages/scheithauer/publications/Weil\_representation.pdf.\par
W. Stein, Modular Forms, a Computational Approach, online Chapter 5, Eisenstein-series normalization and modularity. https://wstein.org/books/modform/modform/eisenstein.html.\par
W. P. Thurston, Shapes of polyhedra and triangulations of the sphere, Geometry \& Topology Monographs 1 (1998), 511–549; selected completion and local-link statements on pp. 512, 528–533. https://msp.org/gtm/1998/01/gtm-1998-01-025p.pdf.\par
\end{document}
